\documentclass{handout}
\usepackage{handoutSetup}\usepackage{handoutShortcuts}  

\begin{document}
\handoutHeader

\begin{verbatimwrite}{\jobname.title}
The Random Walk Model of Consumption
\end{verbatimwrite}

\handoutNameMake

\begin{verbatimwrite}{./body.tex}

This handout derives the \cite{hallRandomWalk} random walk proposition for
consumption.

The consumption Euler equation when future consumption is uncertain
takes the form\footnote{See handout \handoutC{Envelope} for the derivation
of the Euler equation in the perfect foresight case; we will show later 
that the consequence of uncertainty is simply to insert the
expectations operator.}
\begin{equation}\begin{gathered}\begin{aligned}
  \uFunc^{\prime}(\cRat_{t}) & =  \Discount \Rfree \Ex_{t}[\uFunc^{\prime}(\cRat_{t+1})]. \label{eq:uPEuler}
\end{aligned}\end{gathered}\end{equation}

Suppose the utility function is quadratic:
\begin{equation}\begin{gathered}\begin{aligned}
  \uFunc(\cRat) & =  -(1/2) (\cancel{\cRat}-\cRat)^{2}
\end{aligned}\end{gathered}\end{equation}
where $\cancel{\cRat}$ is the ``bliss point'' level of consumption.\footnote{Assume that the consumer
is sufficiently poor that it will be impossible for them ever to achieve consumption as large as $\cancel{\cRat}$.}  Marginal
utility is 
\begin{equation}\begin{gathered}\begin{aligned}
  \uFunc^{\prime}(\cRat) & =   (\cancel{\cRat}-\cRat)
\end{aligned}\end{gathered}\end{equation}
and suppose further that $\Rfree \Discount = 1$ so that \eqref{eq:uPEuler} becomes
\begin{equation}\begin{gathered}\begin{aligned}
  (\cancel{\cRat}-\cRat_{t})  & =  \Ex_{t}[(\cancel{\cRat}-\cRat_{t+1})]
\\ \Ex_{t}[\cRat_{t+1}] & =  \cRat_{t}.
\end{aligned}\end{gathered}\end{equation}

Defining the innovation to consumption as
\begin{equation}\begin{gathered}\begin{aligned}
  \epsilon_{t+1} & =  \cRat_{t+1}-\cRat_{t},
\\ & \equiv  \Delta \cRat_{t+1}
,
\end{aligned}\end{gathered}\end{equation}
the random walk proposition is simply that 
the expectation of consumption changes is zero:
\begin{equation}\begin{gathered}\begin{aligned}
  \Ex_{t}[\Delta \cRat_{t+1}] & =  0.
\end{aligned}\end{gathered}\end{equation}

This means that no information known to the consumer when
the consumption choice $\cRat_{t}$ was made can have any predictive
power for how consumption will change between period $t$ and 
$t+1$ (or for any date beyond $t+1$).

\end{verbatimwrite}
\input ./body.tex

\input handoutBibMake

\end{document}

